\subsection{Sklicevanje in literatura}

Z LaTeXom lahko oblikujemo ne le dokumente, ki so visoke kakovosti za tiskanje. Vendar pa je danes veliko dokumentov prebrati na zaslonu računalnika ali mobilne naprave.
Paket \href{https://ctan.org/pkg/hyperref}{\texttt{hyperref}} omogoča ustvarjanje dokumentov PDF z aktivnimi povezavami oz. \textit{hiperpovezavami}.
To pomeni, da lahko kliknemo na sklic in nas ta pripelje do ustreznega dela dokumenta ali na zunanjo spletno stran.
Paket \texttt{hyperref} omogoča tudi ustvarjanje kazal, ki so interaktivni in omogočajo enostavno navigacijo po dokumentu.

Za zdaj bomo uporabili paket \texttt{hyperref} za sklicevanje na literaturo v našem dokumentu.
Seveda, najprej moramo nalovziti paket v preambuli našega dokumenta:

\begin{verbatim}
\usepackage{hyperref}
\end{verbatim}

\begin{framed}
\textbf{Pazite!}\\
Paket \texttt{hyperref} mora biti naložen kot zadnji paket v preambuli, saj lahko povzroči konflikte z drugimi paketi, če ni naložen na koncu.
\end{framed}

\subsubsection{Sklicevanje}

V dokumentu lahko sklicujemo na različne elemente, kot so slike, tabele, enačbe, razdelki, itd.
Za sklicevanje uporabimo ukaza \texttt{{\textbackslash}label\{\textless oznaka\textgreater \}} in \texttt{{\textbackslash}ref\{\textless oznaka\textgreater \}}.
Ukaz \texttt{{\textbackslash}label\{\textless oznaka\textgreater \}} postavimo takoj za element, na katere se želimo sklicevati. Ukaz \texttt{{\textbackslash}ref\{\textless oznaka\textgreater \}} pa uporabimo na mestu, kjer želimo vstaviti sklic. Argument \texttt{\textless oznaka\textgreater } je poljuben niz znakov, ki ga izberemo sami, vendar mora biti edinstven v celotnem dokumentu.

\begin{example}\label{eg_sklici_1}\end{example}\begin{framed}
\textbf{Nasvet}\\
Oznake je priporočljivo oblikovati tako, da so opisne ampak kratke. Na primer, za razdelek o sklicevanju,\texttt{sec:sklicevanje}, je boljša oznaka kot \texttt{s1} ali \texttt{razdelek1}.

Dobra praksa je tudi, da uporabljamo predpone za različne vrste elementov, kot so \texttt{sec:} za razdelke, \texttt{fig:} za slike, \texttt{tab:} za tabele in \texttt{eq:} za enačbe.
\end{framed}

\begin{framed}
\textbf{Opomba}\\
Ukaz \texttt{{\textbackslash}ref} vrne le številko elementa, na primer številko razdelka ali slike. Potrebno je, da ročno dodamo besedilo, kot je ``Razdelek'', ``Slika'' ali ``Tabela'', pred številko.

Obstajajo paketi, kot je \href{https://ctan.org/pkg/cleveref?lang=en}{\texttt{cleveref}}, ki omogočajo samodejno dodajanje besedila pred številko. Na žalost, paket \texttt{cleveref} ni podpira slovenskega jezika, in čeprav je mogoče definirati nove jezike, je trenutno slovenski jezik težko implementirati zaradi sklanjatev (prim. „glej  Slik\textbf{o} 1'' v „na Slik\textbf{i} 1'').
\end{framed}

Če uporabimo ukaz \texttt{{\textbackslash}ref\{\textless oznaka\textgreater \}} z razdelki ``z zvezdico'' (npr. \texttt{{\textbackslash}section*\{...\}}), ki niso oštevilčeni, bo ukaz vrnil nič ali nepričakovano vrednost, saj ti razdelki nimajo številk.
Zato je priporočljivo, da se na razdelke ``z zvezdico'' ne sklicujemo z ukazom \texttt{{\textbackslash}ref}.
Možna alternativa je uporaba ukaza \texttt{{\textbackslash}nameref\{\textless oznaka\textgreater \}}, ki vrne ime razdelka.

Ukaz \texttt{{\textbackslash}pageref\{\textless oznaka\textgreater \}} vrne številko strani, na kateri se nahaja element z dano oznako.
Ta ukaz lahko uporabimo z razdelki ``z zvezdico'', saj nam ni treba pridobiti številke razdelka, ampak le številko strani.

Za plavajoče elemente, kot so slike in tabele, je priporočljivo, da oznako dodamo znotraj okolja \texttt{figure} ali \texttt{table}, takoj po ukazu \texttt{{\textbackslash}caption\{...\}}.
To zagotavlja, da se oznaka nanaša na pravilno številko slike ali tabele.

\begin{example}\label{eg_sklici_2}\end{example}Označujmo in sklicujmo lahko vse, čemur LaTeX dodeli številko. Na primer, točke v seznamih.

\begin{example}\label{eg_sklici_3}\end{example}\subsubsection{Literartura}

\textbf{Klasični način}

Za vključitev literature v dokument LaTeX, lahko uporabimo okolje \texttt{thebibliography}.
To je le poseben seznam, kjer vsaka referenca dobi svojo oznako z ukazom \texttt{{\textbackslash}bibitem\{\textless oznaka\textgreater \}}.
Nato se na referenco sklicujemo z ukazom \texttt{{\textbackslash}cite\{\textless oznaka\textgreater \}}.

\begin{example}\label{eg_literatura_1}\end{example}Ta način je enostaven in hiter, vendar ima nekaj pomanjkljivosti:

\begin{itemize}
\item Ročno moramo vnašati vse podatke o referencah, kar je lahko zamudno in nagnjeno k napakam.
\item Ni standardiziranega formata za vnos referenc, kar lahko vodi do neenotnosti v slogu citiranja.
\item Enostavne spremembe sloga citiranja so težavne, saj moramo ročno prilagoditi vsako referenco.
\item Ni ločnega/struturnega označevanja, ki bi nakazovalo, kateri del bibliografskega vnosa je avtor, naslov, leto izdaje itd.
\item Sortiranje referenc je ročno in lahko vodi do napak. Urejanje (npr. po letu izdaje namesto po priimku avtorja) je prav tako ročno in zamudno.
\item Na splošno to ni pristop, ki bi bil značilen za LaTeX.
\end{itemize}

Zaradi teh razlogov je priporočljivo uporabiti bolj napredne metode za upravljanje literature v LaTeX dokumentih, kot je BibTeX ali/in BibLaTeX.

\textbf{BibLaTeX način}

BibTex je orodje za upravljanje bibliografskih podatkov, ki deluje skupaj z LaTeXom.
Omogoča nam, da shranimo vse naše reference v ločeno datoteko \texttt{.bib}, kar olajša upravljanje in ponovno uporabo referenc v različnih dokumentih.
BibLaTeX je sodobnejša alternativa BibTeX-u, ki ponuja več funkcionalnosti in prilagodljivosti.

Uporaba BibLaTeX-a vključuje naslednje korake:

\begin{enumerate}
\item Ustvarimo datoteko z bibliografskimi podatki z razširitvijo \texttt{.bib}. V tej datoteki so reference shranjene v posebnem formatu.
\item V preambulo našega LaTeX dokumenta vključimo paket \texttt{biblatex} in določimo slog citiranja ter datoteko z bibliografskimi podatki.
\end{enumerate}

Datoteka \texttt{.bib} je sestavljena iz vnosa, kjer vsak vnos predstavlja eno referenco.
Vsak vnos ima svoj edinstven ključ, ki ga uporabimo za sklicevanje na to referenco v našem dokumentu.
Vsak vnos vsebuje različna polja, kot so avtor, naslov, leto izdaje, itd. in izgledajo nekako takole:

\begin{verbatim}
@<vrsta_Vnosa>{<kjluč>,
<polje1>    = {vrednost1},
<polje2>    = {vrednost2},
<polje3>    = {vrednost3}, 
<poljeN>    = {vrednostN}, 
}
\end{verbatim}

kjer je \texttt{\textless vrsta\_Vnosa\textgreater } lahko \texttt{book}, \texttt{article}, \texttt{inproceedings}, \texttt{misc}, itd., odvisno od vrste reference.
\texttt{\textless ključ\textgreater } je edinstven identifikator za ta vnos, ki ga bomo uporabili v ukazu \texttt{{\textbackslash}cite\{\textless ključ\textgreater \}} v našem LaTeX dokumentu.
\texttt{\textless polje1\textgreater }, \texttt{\textless polje2\textgreater }, itd. so različna polja, ki vsebujejo informacije o referenci, kot so avtor, naslov, leto izdaje, itd. Npr. za knjigo \textit{Harry Potter in kamen modrosti} bi bil vnos videti takole:

\begin{verbatim}
@book{hp1,
author    = {J. K. Rowling},
title     = {Harry Potter in kamen modrosti},
publisher = {Založba EPTA},
year      = {1997}
}
\end{verbatim}

Polja so odvisna od vrste vnosa in slog bibligrafije, ki ga uporabljamo (več o tem kasneje). Za knjige so običajna polja \texttt{author}, \texttt{title}, \texttt{publisher} in \texttt{year}. Za članke v revijah bi lahko uporabili polja, kot so \texttt{journal}, \texttt{volume}, \texttt{number} in \texttt{pages}. Za spletne vire je pogosto uporabno polje \texttt{url} in \texttt{note}.
Celoten seznam vrst vnosov in polj najdete v dokumentaciji paketa \href{https://ctan.org/pkg/biblatex?lang=en}{BibLaTeX} ali tem članku na \href{https://en.wikibooks.org/wiki/LaTeX/Bibliography\_Management\#Entry\_and\_field\_types\_in\_.bib\_files}{wikibooks}.

Dobra novica je, da običajno ni potrebno ročno pisati datoteke \texttt{.bib}, saj obstajajo številna orodja in spletne strani, ki lahko samodejno generirajo vnose v pravilnem formatu.
Nekatera orodja, ki jih lahko uporabite za ustvarjanje in upravljanje datotek \texttt{.bib}, vključujejo:

\begin{itemize}
\item \href{https://scholar.google.com/}{Google Scholar}: omogoča izvoz posameznih referenc v format \texttt{.bib}.
\item \href{https://www.cobiss.si/}{COBISS}: slovenski sistem za iskanje in izposojo knjižničnih gradiv, ki omogoča izvoz bibliografskih podatkov v format \texttt{.bib}.
\item \href{https://zbib.org/}{ZoteroBib}: spletno orodje za hitro ustvarjanje bibliografij, ki podpira izvoz v format \texttt{.bib}.
\item \href{https://mathscinet.ams.org/mathscinet}{MathSciNet}: spletna baza podatkov za matematiko.
\item \href{https://dblp.org/}{DBLP}: baza podatkov za računalništvo.
\end{itemize}

Ko smo pripravili našo bazo podatkov referenc (tj. našo datoteko \texttt{.bib}, recimo \texttt{literatura.bib}), moramo v našem LaTeX dokumentu naložiti paket \texttt{biblatex} in povezati datoteko z bibliografskimi podatki.
To storimo z naslednjimi ukazi v preambuli:

\begin{verbatim}
\usepackage{biblatex} % Naloži paket biblatex
\addbibresource{literatura.bib} % Poveže datoteko z bibliografskimi podatki, recimo literatura.bib
\end{verbatim}

V besedilu našega dokumenta se na reference sklicujemo z ukazom \texttt{{\textbackslash}cite\{\textless ključ\textgreater \}}, kjer je \texttt{\textless ključ\textgreater } ključ vnosa v naši datoteki \texttt{.bib}.
Parameter \texttt{\textless ključ\textgreater } lahko je tudi seznam ključev, ločenih z vejico, če želimo sklicevati na več referenc hkrati.

Na koncu dokumenta (ali, kjerkoli  že želimo natisniti seznam referenc), dodamo ukaz \texttt{{\textbackslash}printbibliography}, ki natisne seznam vseh referenc, na katere smo se sklicevali v besedilu.
Upoštevajte, da čeprav nismo uporabili vseh referenc iz datoteke \texttt{.bib}, bo \texttt{{\textbackslash}printbibliography} natisnil le tiste, na katere smo se sklicevali z ukazom \texttt{{\textbackslash}cite\{\textless ključ\textgreater \}}.

Z uporabo BibLaTeX lahko na nadzorovan način prilagodimo videz naših citatov in sklicev. To storimo z izbiro sloga citiranja, ki ga določimo kot možnost pri nalaganju paketa \texttt{biblatex}.
Na primer,

\begin{verbatim}
\usepackage[style=alphabetic]{biblatex}
\end{verbatim}

naloži slog \texttt{alphabetic}, ki uporablja črke iz imen avtorjev in leto izdaje za označevanje citatov v besedilu. Privzeto je slog \texttt{numeric}, ki uporablja številke za označevanje citatov.
Obstaja veliko različnih slogov citiranja, ki jih lahko uporabimo, odvisno od naših potreb in preferenc.
Obsežen seznam slogov najdete v dokumentaciji paketa \href{https://ctan.org/pkg/biblatex?lang=en}{BibLaTeX} ali na \href{https://www.overleaf.com/learn/latex/Biblatex\_bibliography\_styles}{Overleaf}.

\begin{framed}
\textbf{Pazite!}\\
Slog citatov lahko spreminjamo neodvisno, to pomeni, da lahko spreminjamo, kaj ukaz \texttt{{\textbackslash}cite} natisne v našem besedilu in kako se natisnejo oznake, ko natisnemo bibliografijo, vendar je dobro ohraniti enak slog.
\end{framed}

Če damo možnost \texttt{sorting=ynt} razvrsti bibliografijo po letu izdaje (\texttt{y}), nato po imenu avtorja (\texttt{n}) in nato po naslovu (\texttt{t}). Privzeto je \texttt{sorting=nty}, kar pomeni, da se najprej razvrsti po imenu avtorja, nato po naslovu in nato po letu izdaje.

Celoten seznam možnosti za razvrščanje najdete v dokumentaciji paketa \href{https://ctan.org/pkg/biblatex?lang=en}{\texttt{BibLaTeX}} ali na \href{https://tug.org/biblatex-cheatsheet/biblatex-cheatsheet.pdf}{BibLaTeX cheat sheet}.

Kot smo že omenili, ukaz \texttt{{\textbackslash}printbibliography} natisne seznam vseh referenc, na katere smo se sklicevali v besedilu.
Če želimo natisniti reference, na katere se nismo sklicevali, lahko uporabimo ukaz \texttt{{\textbackslash}nocite\{\textless ključ\textgreater \}}.
Ta ukaz nam omogoča, da vključimo določene vnose v bibliografijo, tudi če se nanje nismo sklicevali v besedilu.
Če želimo vključiti vse vnose iz naše datoteke \texttt{.bib}, lahko uporabimo ukaz \texttt{{\textbackslash}nocite\{*\}}.

Ukaz \texttt{{\textbackslash}printbibliography} lahko sprejme različne možnosti, ki nadzorujejo videz bibliografije.
Na primer, lahko uporabimo možnost \texttt{title}, da spremenimo naslov bibliografije, kot je prikazano spodaj:

\begin{example}\label{eg_literatura_2}\end{example}Referenca lahko organizirate z uporabo filtrov.
Na primer, lahko natisnemo samo knjige z uporabo možnosti \texttt{type=book}:

\begin{example}\label{eg_literatura_3}\end{example}Na voljo so filtri \texttt{type}, \texttt{keywords} in \texttt{category} in lahko uporabimo tudi njihove negacije. Z ukazom \texttt{{\textbackslash}defbibfilter\{\textless ime\_filtra\textgreater \}} lahko definiramo tudi svoje filtre.

\begin{example}\label{eg_literatura_4}\end{example}\begin{example}\label{eg_literatura_5}\end{example}Privzeto se bibliografija ne prikaže v kazalu vsebine. Če želimo, da se prikaže, lahko uporabimo možnost \texttt{heading}.
Različen vrednosti za možnost \texttt{heading} so v  Table~\ref{tab:heading_options}`.

\begin{table}
\centering
\caption[]{Naslov bibliografije v kazalu.}
\label{tab:heading_options}
\begin{tabular}{p{\dimexpr 0.500\linewidth-2\tabcolsep}p{\dimexpr 0.500\linewidth-2\tabcolsep}}
\toprule
Vrednost & Opis \\
\hline
\texttt{bibliography} & privzeto, kot je razdelek >>z zvezdico<<. \\
\texttt{bibintoc} & kot razdelek z zvezdico, vendar se bo pojavil v kazalu. \\
\texttt{subbibliography} & bo bibliografija padla za eno stopnjo v hierarhiji (\texttt{{\textbackslash}section*} namesto \texttt{{\textbackslash}chapter*}, \texttt{{\textbackslash}subsection*} namesto \texttt{{\textbackslash}section*} in tako naprej). \\
\texttt{bibnumbered} & naraven razdelek (brez zvezdice), torej s številko in v kazalu. \\
\texttt{subbibintoc} & podobno kot subbibliography, vendar bo prikazana v kazalu. \\
\texttt{subbibnumbered} & podobno kot subbibliography, vendar brez zvezdice. \\
\texttt{none} & ne bo natisnil naslova. \\
\bottomrule
\end{tabular}
\end{table}

\begin{example}\label{eg_literatura_heading_1}\end{example}\begin{example}\label{eg_literatura_heading_2}\end{example}\begin{example}\label{eg_literatura_heading_3}\end{example}\begin{example}\label{eg_literatura_heading_4}\end{example}\begin{example}\label{eg_literatura_heading_5}\end{example}\begin{example}\label{eg_literatura_heading_6}\end{example}\begin{example}\label{eg_literatura_heading_7}\end{example}\subsubsection{Vaje}\label{latexLiteratura_vaje}

\begin{itemize}
\item Exercise~\ref{ex_hyperref}: Sklicevanje z uporabo paketa \texttt{hyperref}.
\item Exercise~\ref{ex_sklici}: Sklicevanje na različne elemente v dokumentu.
\item Exercise~\ref{ex_biblatex}: Uporaba BibLaTeX-a za upravljanje literature.
\end{itemize}

%  - [ ] {numref}`Vaja {number} <ex_hyperref>`
% - [ ] {numref}`Vaja {number} <ex_sklici>`
% - [ ] {numref}`Vaja {number} <ex_biblatex>`